\section{External Controls and Supplementary Results}

\subsection{Design of Tissue-independent Controls}

The most direct biological control removes tissue identity while retaining peptide and MHC information. A previously reported trainable control of this form was not rerun on the occurrence-matched files, so its legacy performance values are excluded. The revised fixed-test comparison retains only external predictors whose scores were recomputed for every occurrence-matched test row.

Rows are not pooled, deduplicated, or assigned a consensus label after tissue identity is removed. Each original tissue-specific row and label is retained. Consequently, the same peptide--MHC query can carry both labels across tissues: it is positive where ligand evidence was reported and may be a pseudo-negative elsewhere. Table~\ref{tab:tissue-blind-conflict-audit} quantifies this deliberate label ambiguity. A query has one retained label observation when it occurs in one row and multiple retained label observations when it occurs in more than one row; a ``conflicting query'' is a multiple-observation query whose rows contain both binary label values. The conflict percentage therefore uses multiple-observation queries, rather than all unique queries, as its denominator. These are not contradictions within one tissue--MHC task. They arise across tissues because the control is intentionally denied the variable that distinguishes those contexts.

The occurrence-matched control set contains 29,662 unique Human peptide--MHC queries across 29 HLA alleles and 3,338 Mouse queries across four H2 restrictions. We score every test query with MHCflurry and NetMHCpan, two established predictors that do not use target-tissue identity \cite{odonnell2020mhcflurry,reynisson2020netmhcpan}. These fixed external scorers are not retrained on TissuePMHC labels and have no seed-to-seed training variation.

\begin{table*}[t]
\centering
\caption{Deterministic label-conflict audit for controls without tissue information. The table asks how often identical peptide--MHC inputs acquire different tissue-specific labels after tissue identity is removed; it is not a model-performance result and has no training seeds. Queries are defined by peptide sequence and MHC restriction. ``One label'' denotes one retained labeled row; ``multiple labels'' denotes more than one retained labeled row, whether or not their binary values agree. A conflicting query is a multiple-label query containing both 0 and 1. Conflict percentages use multiple-label queries as the denominator. Human and Mouse use the same definitions and partition order.}
\label{tab:tissue-blind-conflict-audit}
\fontsize{7.2}{8.2}\selectfont
\setlength{\tabcolsep}{2.5pt}
\begin{tabularx}{\textwidth}{@{}l>{\raggedright\arraybackslash}Xrrrrrr@{}}
\toprule
Species & Partition & Rows & \shortstack{Unique\\queries} & \shortstack{One retained\\label} & \shortstack{Multiple retained\\labels} & \shortstack{Conflicting\\queries} & \shortstack{Conflict among\\multiple (\%)} \\
\midrule
Human & Training pool & 42,978 & 27,063 & 12,572 & 14,491 & 14,168 & 97.77\% \\
Human & Fixed internal test & 7,700 & 7,012 & 6,327 & 685 & 652 & 95.18\% \\
Human & Complete benchmark & 50,678 & 29,662 & 10,441 & 19,221 & 18,882 & 98.24\% \\
\midrule
Mouse & Training pool & 4,178 & 2,934 & 1,735 & 1,199 & 1,175 & 98.00\% \\
Mouse & Fixed internal test & 1,100 & 996 & 892 & 104 & 103 & 99.04\% \\
Mouse & Complete benchmark & 5,278 & 3,338 & 1,487 & 1,851 & 1,833 & 99.03\% \\
\bottomrule
\end{tabularx}
\end{table*}

Among queries represented by multiple tissue-specific rows, 95.18--99.04\% contain both labels across the six species--partition summaries. The high conflict fractions show why a trainable tissue-blind control is intrinsically ambiguous: identical observed inputs can correspond to different tissue-specific targets. This table is a data audit, not a trainable-control performance claim.

\subsection{External Tissue-agnostic Predictors and Score Combination}

The occurrence-matched fixed test makes tissue-independent external scores approximately random in both species (Table~\ref{tab:rq6-method-comparison}). To integrate the former external-only and combined-score analyses in one comparable table, we also evaluate a fixed, untrained combination rule. Within each tissue--MHC task and each TissuePMHC seed, the TissuePMHC and external scores are separately converted to percentile ranks and then averaged with equal weight. No weight, threshold, or predictor is selected on the fixed test. The combined metrics and their sample standard deviations are calculated across the same three TissuePMHC seeds; the external-only metrics are deterministic. This is a sensitivity analysis of score fusion, not a separately tuned model.

\begin{table*}[t]
\centering
\caption{Unified external-control analysis on the occurrence-matched fixed tests. The table asks whether deterministic tissue-agnostic peptide--MHC scores explain or improve tissue-conditioned ranking. All metrics are task-macro and rows are directly comparable only within species. ``Tissue'' is tuned TissuePMHC alone; ``combined'' is the unweighted mean of within-task percentile ranks from TissuePMHC and the stated external score. External-only values are deterministic, whereas TissuePMHC and combined values use seeds 20260704--20260706; combined entries are means $\pm$ sample standard deviations. $\Delta$ AUROC is combined minus TissuePMHC for the corresponding seed, averaged over seeds. Pair accuracy is the fraction of complete matched pairs with a higher score for the recorded positive. Human and Mouse use the same controls, columns, and display order.}
\label{tab:rq6-method-comparison}
\fontsize{7.0}{8.0}\selectfont
\setlength{\tabcolsep}{1.6pt}
\begin{tabularx}{\textwidth}{@{}l>{\raggedright\arraybackslash}Xrrrrrrr@{}}
\toprule
Species & External control & \shortstack{External\\AUROC} & \shortstack{External\\AUPRC} & \shortstack{Tissue\\AUROC} & \shortstack{Combined\\AUROC} & \shortstack{$\Delta$\\AUROC} & \shortstack{Combined\\AUPRC} & \shortstack{Combined\\pair acc.} \\
\midrule
Human & MHCflurry presentation score & 0.5044 & 0.5299 & 0.8136 & 0.7188 $\pm$ 0.0022 & -0.0948 & 0.7032 $\pm$ 0.0008 & 0.7241 $\pm$ 0.0044 \\
Human & NetMHCpan binding-affinity rank & 0.4874 & 0.5144 & 0.8136 & 0.7034 $\pm$ 0.0016 & -0.1102 & 0.6818 $\pm$ 0.0013 & 0.7057 $\pm$ 0.0022 \\
Human & NetMHCpan eluted-ligand rank & 0.4948 & 0.5193 & 0.8136 & 0.7101 $\pm$ 0.0017 & -0.1035 & 0.6901 $\pm$ 0.0017 & 0.7145 $\pm$ 0.0031 \\
\midrule
Mouse & MHCflurry presentation score & 0.5009 & 0.5266 & 0.8194 & 0.7183 $\pm$ 0.0115 & -0.1011 & 0.7188 $\pm$ 0.0067 & 0.7303 $\pm$ 0.0141 \\
Mouse & NetMHCpan binding-affinity rank & 0.4917 & 0.5249 & 0.8194 & 0.7153 $\pm$ 0.0106 & -0.1042 & 0.7055 $\pm$ 0.0055 & 0.7297 $\pm$ 0.0076 \\
Mouse & NetMHCpan eluted-ligand rank & 0.4924 & 0.5245 & 0.8194 & 0.7169 $\pm$ 0.0104 & -0.1026 & 0.7079 $\pm$ 0.0039 & 0.7339 $\pm$ 0.0115 \\
\bottomrule
\end{tabularx}
\end{table*}

All three external scores remain near random, and equal-weight fusion reduces AUROC by 0.0948--0.1102 in Human and 0.1011--0.1042 in Mouse relative to TissuePMHC alone. Thus, the general peptide--MHC scores neither explain the tissue-conditioned ranking nor improve it under this prespecified fusion rule. The complete architecture comparisons are reported once, in Tables~\ref{tab:human-standard} and \ref{tab:mouse-baselines}, rather than repeated in a supplementary table.
